\chapter{3 次元 Euclid 空間と直交座標系の関係}
    % \section{勾配}
    %     \label{勾配}
    %     \begin{shadebox}
    %         既出の記号の定義は引き継ぐ。
    %         $\phi:\realNumbers^3\to\complexNumbers$は$\realNumbers^3$上で$\mathrm{C}^1$級であるとする。
    %         $\tilde{\phi}(\bm{u}) \coloneq \phi(T(\bm{u}))$とする。
    %         $\phi$の勾配$\nabla_\vr\phi(\vr)$は$\bm{u}$を用いて次式で表せる。
    %         \[ \nabla_\vr\phi(\vr) = \sum_{i=1}^3\frac{\bm{j}_i}{h_i}\partDeriv{\phi}{u_i}{} = \bracks*{\frac{\bm{j}_1}{h_1}, \frac{\bm{j}_2}{h_2}, \frac{\bm{j}_3}{h_3}}^\top\nabla_\bm{u}\tilde{\phi}(\bm{u}) \]
    %     \end{shadebox}
    %     \begin{proof}
    %         \begin{gather*}
    %             \bm{j}_i = \frac{1}{h_i}\partDeriv{\vr}{u_i}{} = \frac{1}{h_i}\partDerivLong{\sum_{j=1}^3 r_j\bm{i}_j}{u_i}{} = \frac{1}{h_i}[\bm{i}_1,\bm{i}_2,\bm{i}_3]\bracks*{\partDeriv{r_1}{u_i}{},\partDeriv{r_2}{u_i}{},\partDeriv{r_3}{u_i}{}}^\top \\
    %             \therefore\;[\bm{j}_1,\bm{j}_2,\bm{j}_3] = [\bm{i}_1,\bm{i}_2,\bm{i}_3]\bracks*{\frac{1}{h_j}\partDeriv{r_i}{u_j}{}}_{(i,j)\in\{1,2,3\}^2} = JH^{-1}
    %         \end{gather*}
    %         ここに$\quad H \coloneq \diag{h_1,h_2,h_3}$であり，$J$は$T$のJacobi行列である。
    %         以上より次式が成り立つ。
    %         \[ [\bm{i}_1,\bm{i}_2,\bm{i}_3] = [\bm{j}_1,\bm{j}_2,\bm{j}_3]HJ^{-1} \]
    %         $\partDeriv{\phi}{r_i}{}$を評価すると次式を得る。
    %         \begin{gather*}
    %             \partDeriv{\phi}{r_i}{} = \sum_{j=1}^3\partDeriv{\tilde{\phi}}{u_j}{}\partDeriv{u_j}{r_i}{} \\
    %             \begin{split}
    %                 \therefore \quad \nabla_\vr\phi(\vr) &= \bracks*{\partDeriv{\phi}{r_1}{}, \partDeriv{\phi}{r_2}{}, \partDeriv{\phi}{r_3}{}}^\top = \bracks*{\partDeriv{u_j}{r_i}{}}_{(i,j)\in\{1,2,3\}^2}\bracks*{\partDeriv{\tilde{\phi}}{u_1}{},\partDeriv{\tilde{\phi}}{u_2}{},\partDeriv{\tilde{\phi}}{u_3}{}}^\top \\
    %                 &= \parens*{J^{-1}}^\top\nabla_\bm{u}\tilde{\phi}(\bm{u})
    %             \end{split}
    %         \end{gather*}
    %         最後の等号は$\bm{u} = T^{-1}(\vr)$と逆関数定理による。
    %         以上より，次式が成り立つ。
    %         \[ \nabla_\vr\phi(\vr) = [\bm{j}_1,\bm{j}_2,\bm{j}_3]HJ^{-1}\parens*{J^{-1}}^\top\nabla_\bm{u}\tilde{\phi}(\bm{u}) \]
    %         $J^{-1}\parens*{J^{-1}}^\top = J^{-1}\parens*{J^\top}^{-1} = \parens*{J^\top J}^{-1}$であり，$T$に関する仮定から$J$の列は互いに直交し，計量係数の定義から$J^\top J$の第$i$対角成分は$h_i^2$である。
    %         よって$\parens*{J^\top J}^{-1} = H^{-2}$であり，次式が成り立つ。
    %         \[ \nabla_\vr\phi(\vr) = [\bm{j}_1,\bm{j}_2,\bm{j}_3]H^{-1}\nabla_\bm{u}\tilde{\phi}(\bm{u}) = \bracks*{\frac{\bm{j}_1}{h_1}, \frac{\bm{j}_2}{h_2}, \frac{\bm{j}_3}{h_3}}^\top\nabla_\bm{u}\tilde{\phi}(\bm{u}) \]
    %     \end{proof}
    % \section{発散}
    %     \label{発散}
    %     \begin{shadebox}
    %         既出の記号の定義は引き継ぐ。
    %         さらに$T$は$\mathrm{C}^2$級であるとし，$\partDeriv{h_i}{u_i}{} = 0\;(i=1,2,3)$が成り立つとする(例: 円柱座標系, 球座標系)。
    %         3次元Euclid空間上の$\mathrm{C}^1$級のベクトル場$\vA(\vr) = \bm{i}_1 A_1(\vr) + \bm{i}_2 A_2(\vr) + \bm{i}_3 A_3(\vr)$を$u_1,u_2,u_3$直交座標系で表したものを$\bm{j}_1 \tilde{A}_1(\bm{u}) + \bm{j}_2 \tilde{A}_2(\bm{u}) + \bm{j}_3 \tilde{A}_3(\bm{u})$とするとき，次式が成り立つ。
    %         \[ \nabla_\vr\cdot\vA(\vr) = \frac{1}{h_1 h_2 h_3}\bracks*{ \partDeriv{h_2 h_3 \tilde{A}_1(\bm{u})}{u_1}{} + \partDeriv{h_3 h_1 \tilde{A}_2(\bm{u})}{u_2}{} + \partDeriv{h_1 h_2 \tilde{A}_3(\bm{u})}{u_3}{} } \]
    %     \end{shadebox}
    %     \begin{proof}
    %         \[ \nabla_\vr\cdot\vA(\vr) = \sum_{i=1}^3 \partDeriv{A_i(\vr)}{r_i}{} = \sum_{i=1}^3 \sum_{j=1}^3 \partDeriv{A_i(\vr)}{u_j}{}\partDeriv{u_j}{r_i}{} = \sum_{i=1}^3 \sum_{j=1}^3 \parens*{\partDerivLong{\sum_{k=1}^3\frac{1}{h_k}\partDeriv{r_i}{u_k}{}\tilde{A}_k(\bm{u})}{u_j}{}}\partDeriv{u_j}{r_i}{} \tag{1} \]
    %         最後の等号成立は\ref{ベクトルの成分の変換}による。
    %         \par
    %         まず$\partDeriv{u_j}{r_i}{}$を評価する。
    %         $J$を$T$のJacobi行列とすると次式が成り立つ。
    %         \[ \partDeriv{u_j}{r_i}{} = J^{-1}[j,i] = \parens*{H^{-2}J^\top}[j,i] = \frac{1}{h_j^2}J[i,j] = \frac{1}{h_j^2}\partDeriv{r_i}{u_j}{} \]
    %         式(1)の$\tilde{A}_k(\bm{u})$に関する項に着目すると，以下のように計算される。
    %         但し以下の式で$\delta_{i,j}$はKroneckerのデルタを表す。
    %         \begin{align*}
    %             \phantom{=} &\sum_{i=1}^3 \sum_{j=1}^3 \parens*{\partDerivLong{\frac{1}{h_k}\partDeriv{r_i}{u_k}{}\tilde{A}_k(\bm{u})}{u_j}{}}\frac{1}{h_j^2}\partDeriv{r_i}{u_j}{} = \sum_{j=1}^3 \frac{1}{h_j^2} \sum_{i=1}^3 \parens*{\partDerivLong{\frac{1}{h_k}\partDeriv{r_i}{u_k}{}\tilde{A}_k(\bm{u})}{u_j}{}}\partDeriv{r_i}{u_j}{} \\
    %             = &\sum_{j=1}^3 \frac{1}{h_j^2} \sum_{i=1}^3 \bracks*{ \parens*{\partDerivLong{\frac{1}{h_k}}{u_j}{}}\tilde{A}_k(\bm{u})\partDeriv{r_i}{u_k}{}\partDeriv{r_i}{u_j}{} + \frac{1}{h_k}\partDeriv{\tilde{A}_k(\bm{u})}{u_j}{}\partDeriv{r_i}{u_k}{}\partDeriv{r_i}{u_j}{} + \frac{\tilde{A}_k(\bm{u})}{h_k}\partDerivIIHetero{r_i}{u_j}{u_k}\partDeriv{r_i}{u_j}{} } \\
    %             = &\sum_{j=1}^3 \frac{1}{h_j^2} \bracks*{ \parens*{\partDerivLong{\frac{1}{h_k}}{u_j}{}}\tilde{A}_k(\bm{u})\partDeriv{\vr}{u_k}{}\cdot\partDeriv{\vr}{u_j}{} + \frac{1}{h_k}\partDeriv{\tilde{A}_k(\bm{u})}{u_j}{}\partDeriv{\vr}{u_k}{}\cdot\partDeriv{\vr}{u_j}{} + \frac{\tilde{A}_k(\bm{u})}{h_k}\partDerivIIHetero{\vr}{u_j}{u_k}\cdot\partDeriv{\vr}{u_j}{} } \\
    %             = &\sum_{j=1}^3 \frac{1}{h_j^2} \bracks*{ \parens*{\partDerivLong{\frac{1}{h_k}}{u_j}{}}\tilde{A}_k(\bm{u})h_k^2\delta_{k,j} + \frac{1}{h_k}\partDeriv{\tilde{A}_k(\bm{u})}{u_j}{}h_k^2\delta_{k,j} + \frac{\tilde{A}_k(\bm{u})}{h_k}\partDerivIIHetero{\vr}{u_j}{u_k}\cdot\partDeriv{\vr}{u_j}{} } \\
    %             = &\frac{1}{h_k^2} \bracks*{ \parens*{\partDerivLong{\frac{1}{h_k}}{u_k}{}}\tilde{A}_k(\bm{u})h_k^2 + h_k\partDeriv{\tilde{A}_k(\bm{u})}{u_k}{} } + \frac{\tilde{A}_k(\bm{u})}{h_k} \sum_{j=1}^3 \frac{1}{h_j^2} \partDerivIIHetero{\vr}{u_j}{u_k}\cdot\partDeriv{\vr}{u_j}{} \\
    %             = &\frac{1}{h_k}\partDeriv{\tilde{A}_k(\bm{u})}{u_k}{} + \frac{\tilde{A}_k(\bm{u})}{h_k} \sum_{j=1}^3 \underbrace{\frac{1}{h_j^2}\partDerivIIHetero{\vr}{u_j}{u_k}\cdot\partDeriv{\vr}{u_j}{}}_{(2)} \quad \parens*{\because \text{仮定より}\partDerivLong{\frac{1}{h_k}}{u_j}{} = 0} \tag{3}
    %         \end{align*}
    %         式(2)を評価する。
    %         $j=k$のとき次式より0である。
    %         \[ \partDeriv{\vr}{u_k}{2}\cdot\partDeriv{\vr}{u_k}{} = \frac{1}{2}\partDerivLong{\parens*{\partDeriv{\vr}{u_k}{}\cdot\partDeriv{\vr}{u_k}{}}}{u_k}{} = \frac{1}{2}\partDeriv{h_k^2}{u_k}{} = 0 \]
    %         $j\neq k$のとき次式を得る。
    %         \begin{align*}
    %             \phantom{=} &\frac{1}{h_j^2}\partDerivIIHetero{\vr}{u_j}{u_k}\cdot\partDeriv{\vr}{u_j}{} = \frac{1}{h_j^2}\partDerivIIHetero{\vr}{u_k}{u_j}\cdot\partDeriv{\vr}{u_j}{} \\
    %             &= \frac{1}{h_j^2}\times\frac{1}{2} \partDerivLong{\parens*{\partDeriv{\vr}{u_j}{}\cdot\partDeriv{\vr}{u_j}{}}}{u_k}{} = \frac{1}{h_j^2}\times\frac{1}{2}\partDeriv{h_j^2}{u_k}{} = \frac{1}{h_j}\partDeriv{h_j}{u_k}{}
    %         \end{align*}
    %         これらを式(3)に適用して次式を得る。
    %         但し以下で$[i+j]$は$i+j$を3で割った余りを表す。
    %         \[ (3) = \frac{1}{h_k}\partDeriv{\tilde{A}_k(\bm{u})}{u_k}{} + \frac{\tilde{A}_k(\bm{u})}{h_k}\bracks*{ \frac{1}{h_{[k+1]}}\partDeriv{h_{[k+1]}}{u_k}{} + \frac{1}{h_{[k+2]}}\partDeriv{h_{[k+2]}}{u_k}{} } = \frac{1}{h_1 h_2 h_3}\partDeriv{h_{[k+1]}h_{[k+2]}\tilde{A}_k(\bm{u})}{u_k}{} \]
    %     \end{proof}
    % \section{回転}
    %     \begin{shadebox}
    %         既出の記号の定義は引き継ぐ。
    %         さらに$T$は$\mathrm{C}^2$級であるとし，次式が成り立つとする(例: 円柱座標系, 球座標系)。
    %         \[ \partDerivIIHetero{\vr}{u_j}{u_i}\cdot\partDeriv{\vr}{u_k}{} = 0 \quad \parens*{\{i,j,k\} = \{1,2,3\}} \tag{1} \]
    %         3次元Euclid空間上の$\mathrm{C}^1$級のベクトル場$\vA(\vr) = \bm{i}_1 A_1(\vr) + \bm{i}_2 A_2(\vr) + \bm{i}_3 A_3(\vr)$を$u_1,u_2,u_3$直交座標系で表したものを$\bm{j}_1 \tilde{A}_1(\bm{u}) + \bm{j}_2 \tilde{A}_2(\bm{u}) + \bm{j}_3 \tilde{A}_3(\bm{u})$とするとき，次式が成り立つ。
    %         \begin{align*}
    %             \nabla_\vr\times\vA(\vr) &= \frac{\bm{j}_1}{h_2 h_3}\bracks*{ \partDeriv{h_3\tilde{A}_3}{u_2}{} - \partDeriv{h_2\tilde{A}_2}{u_3}{} } + \frac{\bm{j}_2}{h_3 h_1}\bracks*{ \partDeriv{h_1\tilde{A}_1}{u_3}{} - \partDeriv{h_3\tilde{A}_3}{u_1}{} } + \frac{\bm{j}_3}{h_1 h_2}\bracks*{ \partDeriv{h_2\tilde{A}_2}{u_1}{} - \partDeriv{h_1\tilde{A}_1}{u_2}{} } \\
    %             &= \sum_{i=1}^3 \bm{j}_i \varepsilon_{ijk}\frac{1}{h_j h_k}\partDeriv{h_k \tilde{A}_k}{u_j}{}
    %             % = &\sum_{i=1}^3 \frac{\bm{j}_i}{h_{[i+1]}h_{[i+2]}}\bracks*{ \partDeriv{h_{[i+2]}\tilde{A}_{[i+2]}}{u_{[i+1]}}{} - \partDeriv{h_{[i+1]}\tilde{A}_{[i+1]}}{u_{[i+2]}}{} }
    %         \end{align*}
    %         % ここに$[i+j]$は$i+j$を3で割った余りを表す。
    %     \end{shadebox}
    %     \begin{proof}
    %         \begin{align*}
    %             \phantom{=} &\nabla_\vr\times\vA(\vr) = \nabla_\vr \times \sum_{i=1}^3 A_i(\vr)\bm{i}_i = \sum_{i=1}^3 \nabla_\vr \times A_i(\vr)\bm{i}_i = \sum_{i=1}^3 (\nabla_\vr A_i(\vr))\times\bm{i}_i\quad (\because \text{\ref{ベクトルのスカラー関数倍の回転}}) \\
    %             = &\sum_{i=1}^3 \sum_{j=1}^3 \frac{\bm{j}_j}{h_j}\partDeriv{A_i(\vr)}{u_j}{}\times\bm{i}_i = \sum_{i=1}^3 \sum_{j=1}^3 \frac{\bm{j}_j}{h_j}\parens*{\partDerivLong{\sum_{k=1}^3 \frac{\tilde{A}_k(\bm{u})}{h_k}\partDeriv{r_i}{u_k}{}}{u_j}{}} \times \sum_{l=1}^3 \frac{\bm{j}_l}{h_l}\partDeriv{r_i}{u_l}{} \\
    %             &(\because \text{\ref{勾配}, \ref{ベクトルの成分の変換}, \ref{単位ベクトルの変換}})
    %         \end{align*}
    %         以下では表記の簡便さのため，$\tilde{A}_k(\bm{u})$を単に$\tilde{A}_k$と書く。
    %         上式の第$m\in\{1,2,3\}$成分を評価する。
    %         第$m$成分に寄与するのは$(j,l)=\parens*{[m+1],[m+2]},\parens*{[m+2],[m+1]}$のときであり，計算すると次式を得る。
    %         \begin{align*}
    %             \phantom{=} &\sum_{i=1}^3 \left[ \frac{1}{h_{[m+1]}}\parens*{\partDerivLong{\sum_{k=1}^3 \frac{\tilde{A}_k(\bm{u})}{h_k}\partDeriv{r_i}{u_k}{}}{u_{[m+1]}}{}}\frac{1}{h_{[m+2]}}\partDeriv{r_i}{u_{[m+2]}}{} \right. \\
    %             &\left. - \frac{1}{h_{[m+2]}}\parens*{\partDerivLong{\sum_{k=1}^3 \frac{\tilde{A}_k(\bm{u})}{h_k}\partDeriv{r_i}{u_k}{}}{u_{[m+2]}}{}}\frac{1}{h_{[m+1]}}\partDeriv{r_i}{u_{[m+1]}}{} \right] \\
    %             &= \frac{1}{h_{[m+1]}h_{[m+2]}}\sum_{k=1}^3\bracks*{ \sum_{i=1}^3\parens*{\partDerivLong{\frac{\tilde{A}_k}{h_k}\partDeriv{r_i}{u_k}{}}{u_{[m+1]}}{}}\partDeriv{r_i}{u_{[m+2]}}{} - \sum_{i=1}^3\parens*{\partDerivLong{\frac{\tilde{A}_k}{h_k}\partDeriv{r_i}{u_k}{}}{u_{[m+2]}}{}}\partDeriv{r_i}{u_{[m+1]}}{} } \tag{2}
    %         \end{align*}
    %         この式の[ ]内第1項を評価すると次式を得る。
    %         \begin{align*}
    %             \phantom{=} &\sum_{k=1}^3\sum_{i=1}^3\parens*{\partDerivLong{\frac{\tilde{A}_k}{h_k}\partDeriv{r_i}{u_k}{}}{u_{[m+1]}}{}}\partDeriv{r_i}{u_{[m+2]}}{} \\
    %             = &\sum_{k=1}^3 \bracks*{ \parens*{\partDerivLong{\frac{\tilde{A}_k}{h_k}}{u_{[m+1]}}{}}\sum_{i=1}^3\partDeriv{r_i}{u_k}{}\partDeriv{r_i}{u_{[m+2]}}{} + \frac{\tilde{A}_k}{h_k}\sum_{i=1}^3\partDerivIIHetero{r_i}{u_{[m+1]}}{u_k}\partDeriv{r_i}{u_{[m+2]}}{} } \\
    %             = &\sum_{k=1}^3 \bracks*{ \parens*{\partDerivLong{\frac{\tilde{A}_k}{h_k}}{u_{[m+1]}}{}}\partDeriv{\vr}{u_k}{}\cdot\partDeriv{\vr}{u_{[m+2]}}{} + \frac{\tilde{A}_k}{h_k}\partDerivIIHetero{\vr}{u_{[m+1]}}{u_k}\cdot\partDeriv{\vr}{u_{[m+2]}}{} } \\
    %             = &\parens*{\partDerivLong{\frac{\tilde{A}_{[m+2]}}{h_{[m+2]}}}{u_{[m+1]}}{}}h_{[m+2]}^2 + \sum_{k\in\{[m+1],[m+2]\}} \frac{\tilde{A}_k}{h_k}\partDerivIIHetero{\vr}{u_{[m+1]}}{u_k}\cdot\partDeriv{\vr}{u_{[m+2]}}{}
    %         \end{align*}
    %         最後の等号成立には式(1)を用いた。
    %         式(2)の[ ]内第2項も同様に評価すると，式(2)の$1/(h_{[m+1]}h_{[m+2]})$の右側の部分は次式になる。
    %         \begin{align*}
    %             \phantom{=} &\parens*{\partDerivLong{\frac{\tilde{A}_{[m+2]}}{h_{[m+2]}}}{u_{[m+1]}}{}}h_{[m+2]}^2 - \parens*{\partDerivLong{\frac{\tilde{A}_{[m+1]}}{h_{[m+1]}}}{u_{[m+2]}}{}}h_{[m+1]}^2 \\
    %             \phantom{=}    & + \sum_{k\in\{[m+1],[m+2]\}}\frac{\tilde{A}_k}{h_k}\parens*{\partDerivIIHetero{\vr}{u_{[m+1]}}{u_k}\cdot\partDeriv{\vr}{u_{[m+2]}}{} - \partDerivIIHetero{\vr}{u_{[m+2]}}{u_k}\cdot\partDeriv{\vr}{u_{[m+1]}}{}} \tag{3}
    %         \end{align*}
    %         この式の第3項の()内を評価する。
    %         $T$は$\mathrm{C}^2$級であると仮定しているから，微分順序の入れ替えが可能であることに留意する。
    %         $k=[m+1]$のときは次式になる。
    %         \begin{align*}
    %             \phantom{=} &\partDeriv{\vr}{u_{[m+1]}}{2}\cdot\partDeriv{\vr}{u_{[m+2]}}{} - \partDerivIIHetero{\vr}{u_{[m+2]}}{u_{[m+1]}}\cdot\partDeriv{\vr}{u_{[m+1]}}{} \\
    %             = &\partDerivLong{\parens*{\partDeriv{\vr}{u_{[m+1]}}{}\cdot\partDeriv{\vr}{u_{[m+2]}}{}}}{u_{[m+1]}}{} - \partDeriv{\vr}{u_{[m+1]}}{}\cdot\partDerivIIHetero{\vr}{u_{[m+1]}}{u_{[m+2]}}{} - \partDerivIIHetero{\vr}{u_{[m+2]}}{u_{[m+1]}}\cdot\partDeriv{\vr}{u_{[m+1]}}{} \\
    %             = &-\partDeriv{\vr}{u_{[m+1]}}{}\cdot\partDerivIIHetero{\vr}{u_{[m+2]}}{u_{[m+1]}}{} - \partDerivIIHetero{\vr}{u_{[m+2]}}{u_{[m+1]}}\cdot\partDeriv{\vr}{u_{[m+1]}}{} \\
    %             = &-\partDerivLong{\parens*{\partDeriv{\vr}{u_{[m+1]}}{}\cdot\partDeriv{\vr}{u_{[m+1]}}{}}}{u_{[m+2]}}{} = -\partDeriv{h_{[m+1]}^2}{u_{[m+2]}}{}
    %         \end{align*}
    %         $k=[m+2]$のときは次式になる。
    %         \begin{align*}
    %             \phantom{=} &\partDerivIIHetero{\vr}{u_{[m+1]}}{u_{[m+2]}}\cdot\partDeriv{\vr}{u_{[m+2]}}{} - \partDeriv{\vr}{u_{[m+2]}}{2}\cdot\partDeriv{\vr}{u_{[m+1]}}{} \\
    %             = &\partDerivIIHetero{\vr}{u_{[m+1]}}{u_{[m+2]}}\cdot\partDeriv{\vr}{u_{[m+2]}}{} - \partDerivLong{\parens*{\partDeriv{\vr}{u_{[m+2]}}{}\cdot\partDeriv{\vr}{u_{[m+1]}}{}}}{u_{[m+2]}}{} + \partDeriv{\vr}{u_{[m+2]}}{}\cdot\partDerivIIHetero{\vr}{u_{[m+2]}}{u_{[m+1]}} \\
    %             = &\partDerivIIHetero{\vr}{u_{[m+1]}}{u_{[m+2]}}\cdot\partDeriv{\vr}{u_{[m+2]}}{} + \partDeriv{\vr}{u_{[m+2]}}{}\cdot\partDerivIIHetero{\vr}{u_{[m+1]}}{u_{[m+2]}} \\
    %             = &\partDerivLong{\parens*{\partDeriv{\vr}{u_{[m+2]}}{}\cdot\partDeriv{\vr}{u_{[m+2]}}{}}}{u_{[m+1]}}{} = \partDeriv{h_{[m+2]}^2}{u_{[m+1]}}{}
    %         \end{align*}
    %         これらの式を(3)に適用すると次式を得る。
    %         \begin{align*}
    %             (3) &= \parens*{\partDerivLong{\frac{\tilde{A}_{[m+2]}}{h_{[m+2]}}}{u_{[m+1]}}{}}h_{[m+2]}^2 + \frac{\tilde{A}_{[m+2]}}{h_{[m+2]}}\partDeriv{h_{[m+2]}^2}{u_{[m+1]}}{} - \parens*{\partDerivLong{\frac{\tilde{A}_{[m+1]}}{h_{[m+1]}}}{u_{[m+2]}}{}}h_{[m+1]}^2 - \frac{\tilde{A}_{[m+1]}}{h_{[m+1]}}\partDeriv{h_{[m+1]}^2}{u_{[m+2]}}{} \\
    %             &= \partDerivLong{\frac{\tilde{A}_{[m+2]}}{h_{[m+2]}}h_{[m+2]}^2}{u_{[m+1]}}{} - \partDerivLong{\frac{\tilde{A}_{[m+1]}}{h_{[m+1]}}h_{[m+1]}^2}{u_{[m+2]}}{} \\
    %             &= \partDeriv{h_{[m+2]}\tilde{A}_{[m+2]}}{u_{[m+1]}}{} - \partDeriv{h_{[m+1]}\tilde{A}_{[m+1]}}{u_{[m+2]}}{}
    %         \end{align*}
    %         これを式(2)に適用して定理の主張を得る。
    %     \end{proof}
    % \section{Laplacian}
    %     \begin{shadebox}
    %         既出の記号の定義は引き継ぐ。
    %         さらに$T$は$\mathrm{C}^2$級であるとし，$\partDeriv{h_i}{u_i}{} = 0\;(i=1,2,3)$が成り立つとする(例: 円柱座標系, 球座標系)。
    %         $\phi:\realNumbers^3\to\complexNumbers$は$\realNumbers^3$上で$\mathrm{C}^2$級であるとする。
    %         $\tilde{\phi}(\bm{u}) \coloneq \phi(T(\bm{u}))$とする。
    %         $\nabla^2_\vr\phi(\vr)$は$\bm{u}$を用いて次式で表せる。
    %         \[ \nabla^2_\vr\phi(\vr) = \frac{1}{h_1 h_2 h_3}\bracks*{ \partDerivLong{\parens*{\frac{h_2 h_3}{h_1}\partDeriv{\tilde{\phi}(\bm{u})}{u_1}{}}}{u_1}{} + \partDerivLong{\parens*{\frac{h_3 h_1}{h_2}\partDeriv{\tilde{\phi}(\bm{u})}{u_2}{}}}{u_2}{} + \partDerivLong{\parens*{\frac{h_1 h_2}{h_3}\partDeriv{\tilde{\phi}(\bm{u})}{u_3}{}}}{u_3}{} } \]
    %     \end{shadebox}
    %     \begin{proof}
    %         \ref{勾配},\ref{発散}を用いて容易に示せる。
    %     \end{proof}